\documentclass[aspectratio=169,9pt]{beamer}

\usepackage{fontspec}
\usepackage{polyglossia}
\setmainlanguage{russian}
\setotherlanguage{english}
\tracinglostchars=3

\setmainfont{Noto Serif}
\setsansfont{Noto Sans}
\setmonofont{Noto Sans Mono}

\usepackage{mathtools}
\usepackage[labelfont=bf, labelsep=period]{caption}
\usepackage{booktabs}
\usepackage{colortbl}
\usepackage{graphicx}
\graphicspath{{images/}}

% Тема и цвета.
\usetheme{Berlin}
\usecolortheme{default}

% Пользовательские цвета.
\definecolor{good}{HTML}{1E8449}
\definecolor{amber}{HTML}{B8860B}
\definecolor{accent}{HTML}{C0392B}

% Убираем навигационные символы.
\setbeamertemplate{navigation symbols}{}
\setbeamertemplate{frametitle continuation}{}

% Хедер/футер.
\setbeamertemplate{headline}{}
\setbeamertemplate{footline}{
  \hfill%
  \usebeamercolor[fg]{page number in head/foot}%
  \usebeamerfont{page number in head/foot}%
  \insertframenumber~/ \inserttotalframenumber%
  \hspace{1em}%
  \vspace{0.5em}%
}

\title{Разработка средств автоматического\\отображения между таксономиями}
\subtitle{Embedding-based \enspace\textbullet\enspace LLM (LLMs4OM) \enspace\textbullet\enspace GraphSAGE \enspace\textbullet\enspace OAEI 2025}
\author{Михнович Иван Вадимович}
\institute{НИУ ВШЭ \enspace\textbullet\enspace МОАД \enspace\textbullet\enspace 1-й курс}
\date{Санкт-Петербург, июнь 2026 г.}

\begin{document}

% --- 1. Титульный слайд ---
{
\setbeamertemplate{footline}{}
\begin{frame}[plain]
  \titlepage
\end{frame}
}

% --- 2. Постановка задачи ---
\begin{frame}{Постановка задачи}
  \begin{columns}[T,onlytextwidth]
    \column{0.45\textwidth}
      \structure{\large Таксономии}
      \vspace{0.3em}

      \begin{itemize}
        \item Т.н. <<дерево классификатора>>
        \item Узлы~--- классы (например, \texttt{Document}, \texttt{Paper})
        \item Рёбра~--- отношение «is-a» (\texttt{subClassOf})
        \item В общем случае DAG, а не дерево
        \item Можно рассматривать как частный случай онтологии
      \end{itemize}

    \column{0.52\textwidth}
      \structure{\large Сопоставление таксономий}
      \begin{itemize}
        \item Даны две таксономии \texttt{A} и \texttt{B}
        \item Нужно найти соответствия между их классами
        \item Один и тот же концепт~--- разные имена, разная иерархия
        \item Результат: alignment = список матчей (1:1)
      \end{itemize}

      \vspace{0.3em}
      \structure{\large Сложность}
      \begin{itemize}
        \item Синонимы: \texttt{Person} $\leftrightarrow$ \texttt{Human}
        \item Разная гранулярность: \texttt{Paper} $\leftrightarrow$ \texttt{FullPaper}?
        \item Отсутствие общей лексики в текстовых описаниях узлов ~--- нужно понимать смысл написанного
      \end{itemize}
  \end{columns}
\end{frame}

% --- 3. Пример мэтчинга (диаграмма) ---
{
\setbeamertemplate{footline}{}
\begin{frame}[plain]
  \begin{center}
    \includegraphics[width=0.80\textwidth,height=0.85\textheight,keepaspectratio]{matching_diagram.png}
  \end{center}
\end{frame}
}

% --- 4. Данные: OAEI 2025 Conference Track ---
\begin{frame}{Данные: OAEI 2025 Conference Track}
  \begin{columns}[T,onlytextwidth]
    \column{0.48\textwidth}
      \structure{\large OAEI (Ontology Alignment Evaluation Initiative)}
      \begin{itemize}
        \item Ежегодное соревнование по ontology matching
        \item Conference Track~--- 7 онтологий про научные конференции
        \item 21 референсный alignment (ground truth)
        \item Онтологии были урезаны до таксономий путём уничтожения всех связей кроме \texttt{subClassOf}
        \item 29--140 узлов в каждой онтологии
      \end{itemize}

    \column{0.48\textwidth}
      \structure{\large Synthetic data}
      \begin{itemize}
        \item 4 сценария на базе cmt: exact, synonym, structural, attribute
        \item Использовались для GNN-аугментации
      \end{itemize}
  \end{columns}
\end{frame}

% --- 5. Список онтологий
\begin{frame}{Данные: OAEI 2025 Conference Track (продолжение)}
  \structure{\large Список онтологий}
  \vspace{0.3em}

  \begin{tabular}{@{}ll@{}r@{}}
    \toprule
    Код & Название & Кол-во узлов \\
    \midrule
    cmt & Conference Management Tool & 29 \\
    confOf & Conference Offerings & 38 \\
    conference & Conference & 59 \\
    edas & EDAS & 103 \\
    ekaw & EKAW & 73 \\
    iasted & IASTED & 140 \\
    sigkdd & SIGKDD & 49 \\
    \bottomrule
  \end{tabular}

  \vspace{0.5em}
  \alert{Разный размер, разная лексика, разная структура.}
\end{frame}

% --- 6. Метрики ---
\begin{frame}{Метрики оценки}
  \begin{columns}[T,onlytextwidth]
    \column{0.48\textwidth}
      \structure{\large Основная метрика: F1}
      \vspace{0.3em}

      \begin{itemize}
        \item 1:1 matching с порогом уверенности
        \item Precision = TP / (TP + FP)
        \item Recall = TP / (TP + FN)
        \item F1 = гармоническое среднее
      \end{itemize}

      \vspace{0.5em}
      \structure{\large Дополнительные метрики}
      \begin{itemize}
        \item Precision и Recall отдельно
        \item True Positives (количество правильных пар)
        \item Predicted count vs. Ground Truth count
      \end{itemize}

    \column{0.48\textwidth}
      \structure{\large Baseline: StringEquiv}
      \vspace{0.3em}

      \begin{itemize}
        \item Точное совпадение имён классов без учёта регистра букв
        \item Официальный lower bound OAEI
        \item Если имена совпадают~--- матч
        \item На некоторых парах онтологий достигает F1 = 1{,}0
        \item Ломается на синонимах и любых различиях в лексике
      \end{itemize}

      \vspace{0.5em}
      \begin{center}
        \color{amber}\large
        Любой разумный подход к этой задаче должен быть не хуже, чем StringEquiv
      \end{center}
  \end{columns}
\end{frame}

% --- 7. Подходы (обзор) ---
\begin{frame}{Подходы}
  \begin{enumerate}
    \item \structure{Embedding-based}

      {\small BERT (LaBSE, MiniLM, ruRoberta, rubert) + cosine similarity.\\
      Результаты сильно зависят от того, какой текст мы извлекаем из узлов.\\
      Оказалось что \texttt{name\_only} на OAEI работает лучше, чем имя + путь.}

    \vspace{0.5em}
    \item \structure{LLM-based (LLMs4OM)}

      {\small DeepSeek-v4-flash / GPT-4.1-mini.  Бинарная классификация
      пар-кандидатов при помощи LLM (yes/no).  Три представления: C, CP, CC.}

    \vspace{0.5em}
    \item \structure{GraphSAGE (GNN)}

      {\small Siamese GraphSAGE на MiniLM-фичах.  Leave-one-out CV.\\
      \alert{Отрицательный результат}: хуже чем исходные эмбеддинги сами по себе.}

    \vspace{0.5em}
    \item \structure{Ablation study: BM25 + LLM}

      {\small Лексический ретривер (без и в комбинации с эмбеддингами).\\
      \alert{Гибридный ретривер даёт лучший средний F1} (0{,}685 против 0{,}644 у чистого MiniLM).}
  \end{enumerate}
\end{frame}

% --- 8. Embedding-based подход ---
\begin{frame}{Embedding-based: BERT + cosine similarity}
  \begin{columns}[T,onlytextwidth]
    \column{0.48\textwidth}
      \structure{\large Схема}
      \begin{enumerate}
        \item Имя класса $\to$ sentence-transformer $\to$ эмбеддинг
        \item Для каждого source-узла: cosine similarity со всеми target-узлами
        \item 1:1 greedy matching по убыванию сходства
        \item 4 модели: LaBSE, MiniLM, ruRoberta-large, rubert-base
      \end{enumerate}

      \vspace{0.5em}
      \structure{\large Интересное наблюдение}
      \begin{itemize}
        \item Включение пути/атрибутов в описание узла \alert{ухудшает} F1 в 1,5 раза
        \item Только имя: 0{,}314 | Имя + путь + атрибуты: 0{,}207 (MiniLM)
        \item BERT запутывается в разных иерархиях для одного концепта
      \end{itemize}

    \column{0.48\textwidth}
      \structure{\large Результаты (21 пара)}
      \vspace{0.2em}

      \begin{tabular}{@{}lcc@{}}
        \toprule
        Модель & Dim & Avg F1 \\
        \midrule
        LaBSE & 768 & \textbf{0{,}323} \\
        MiniLM & 384 & 0{,}314 \\
        ruRoberta-large & 1024 & 0{,}279 \\
        rubert-base & 768 & 0{,}258 \\
        \bottomrule
      \end{tabular}

      \vspace{0.5em}
      \textbf{Выводы:}
      \begin{itemize}
        \item LaBSE лучший (multilingual, высокий dim)
        \item MiniLM почти не уступает (вдвое меньше dim, быстрее)
        \item Русскоязычные модели хуже работают на английском тексте (duh!)
        \item Структурная информация запутывает BERTы
      \end{itemize}
  \end{columns}
\end{frame}

% --- 9. LLM-based подход ---
\begin{frame}{LLM-based: LLMs4OM framework}
  \begin{columns}[T,onlytextwidth]
    \column{0.48\textwidth}
      \structure{\large Схема LLMs4OM (arXiv 2404.10317)}
      \begin{enumerate}
        \item MiniLM эмбеддинг $\to$ top-5 кандидатов
        \item Для каждого узла LLM классифицирует каждую потенциальную пару (yes/no)
        \item Три представления для LLM: C (имя), CP (имя + родители), CC (имя + дети)
        \item Жадный отбор пар по убыванию уверенности, threshold $>0{,}7$
        \item Structured JSON output
        \item 10 workers, exponential backoff retry
      \end{enumerate}

      \vspace{0.5em}
      \structure{\large Модели}
      \begin{itemize}
        \item \textbf{GPT-4.1-mini} (OpenRouter)
        \item \textbf{DeepSeek-v4-flash} (DeepSeek Platform)
      \end{itemize}

    \column{0.48\textwidth}
      \structure{\large Пример: cmt $\leftrightarrow$ confOf (DeepSeek)}
      \vspace{0.2em}

      \begin{tabular}{@{}lc@{}}
        \toprule
        Режим & F1 \\
        \midrule
        Embedding (MiniLM) & 0{,}308 \\
        StringEquiv & 0{,}533 \\
        \midrule
        LLM C (имя) & 0{,}667 \\
        LLM CP (+родители) & \textbf{0{,}750} \\
        LLM CC (+дети) & 0{,}588 \\
        \bottomrule
      \end{tabular}

      \vspace{0.5em}
      \textbf{Прирост:} +0{,}217 над StringEquiv (+41\%).

      \vspace{0.3em}
      \textbf{Стоимость:} source\_nodes $\times$ top\_k $\times$ 3\_modes = 29 $\times$ 5 $\times$ 3 = 435 обращений к LLM.
  \end{columns}
\end{frame}

% --- 10. Главный результат: LLM vs baselines (график) ---
\begin{frame}[plain]
  \begin{center}
    \includegraphics[width=0.92\textwidth,height=0.90\textheight,keepaspectratio]{comparison_six_pairs.png}
  \end{center}
\end{frame}

% --- 11. LLM vs StringEquiv: анализ ---
\begin{frame}{LLM vs StringEquiv: анализ}
  \structure{\large LLM превосходит StringEquiv на <<сложных>> парах таксономий}

  \vspace{0.5em}
  \begin{tabular}{@{}lccc@{}}
    \toprule
    Пара & StringEquiv & Лучший результат LLM & $\Delta$ \\
    \midrule
    cmt $\leftrightarrow$ conference & 0{,}375 & 0{,}640 (GPT) & \textbf{+71\%} \\
    cmt $\leftrightarrow$ confOf & 0{,}533 & 0{,}750 (DS) & \textbf{+41\%} \\
    confOf $\leftrightarrow$ conference & 0{,}737 & 0{,}762 (DS) & +3\% \\
    cmt $\leftrightarrow$ edas & \textbf{1{,}000} & 1{,}000 (BM25) & 0\% \\
    confOf $\leftrightarrow$ ekaw & 0{,}552 & 0{,}667 (BM25) & \textbf{+21\%} \\
    conference $\leftrightarrow$ ekaw & 0{,}457 & 0{,}577 (GPT) & \textbf{+26\%} \\
    \bottomrule
  \end{tabular}

  \vspace{0.5em}
  \begin{itemize}
    \item На парах с низким лексическим перекрытием (cmt$\leftrightarrow$conference, 37{,}5\%) LLM показывают лучшие результаты: +71\%.
    \item cmt$\leftrightarrow$edas~--- потолок: все 8 GT-совпадений~--- точное совпадение по строкам, улучшать нечего.
    \item confOf$\leftrightarrow$conference: StringEquiv уже силён (0{,}737), LLM приносит скромный буст +3\%.
    \item \textbf{LLM требуют много времени и ресурсов, но их не нужно дообучать под задачу; выигрыш над StringEquiv~--- только на парах с низким лексическим перекрытием, в среднем силы сопоставимы (0{,}644 у LLM+MiniLM против 0{,}663 у StringEquiv).}
  \end{itemize}
\end{frame}

% --- 7. BM25 ablation ---
\begin{frame}{Ablation study: BM25 и гибридный ретривер}
  \begin{columns}[T,onlytextwidth]
    \column{0.50\textwidth}
      \structure{\large Что проверили}
      \begin{itemize}
        \item \textbf{BM25}: лексический ретривер + DeepSeek
        \item \textbf{Hybrid}: union top-K (MiniLM + BM25) + DeepSeek
      \end{itemize}

      \vspace{0.5em}
      \structure{\large Результат}
      \begin{itemize}
        \item BM25 лучший на 8/21 парах: cmt$\leftrightarrow$ekaw (0{,}857), confOf$\leftrightarrow$sigkdd (0{,}909) и др.
        \item Hybrid лучший на 5/21 парах: edas$\leftrightarrow$ekaw (0{,}811), ekaw$\leftrightarrow$sigkdd (0{,}857) и др.
        \item Оба ретривера улучшают средний F1: 0{,}662 (BM25) и 0{,}685 (Hybrid) против 0{,}644 у чистого MiniLM; гибридный~--- лучший в среднем
      \end{itemize}

    \column{0.46\textwidth}
      \structure{\large Интерпретация}
      \begin{itemize}
        \item BM25 полезен, когда имена классов пересекаются лексически (confOf$\leftrightarrow$sigkdd: \texttt{Activity}, \texttt{Conference})
        \item Hybrid помогает, когда MiniLM-ретривер пропускает хороших кандидатов (edas$\leftrightarrow$ekaw)
        \item На 13 парах из 21 BM25-ретривер лучше чистого MiniLM (Hybrid~--- на 11), так что одного MiniLM недостаточно
      \end{itemize}

      \vspace{0.5em}
      \begin{center}
        \color{amber}\large
        BM25 и Hybrid улучшают средний F1;\\
        лучший средний результат~--- у гибридного ретривера
      \end{center}
  \end{columns}
\end{frame}

% --- 8. GNN: архитектура ---
\begin{frame}[plain]
  \begin{center}
    \includegraphics[width=0.90\textwidth,height=0.88\textheight,keepaspectratio]{gnn_architecture.png}
  \end{center}
\end{frame}

% --- 9. GNN: обучение и результат ---
\begin{frame}{GraphSAGE: как обучали и что получилось}
  \begin{columns}[T,onlytextwidth]
    \column{0.52\textwidth}
      \structure{\large Идея}
      \begin{itemize}
        \item MiniLM знает \textbf{имя} класса (\texttt{Paper})
        \item GNN добавляет знание о \textbf{соседях} (\texttt{Document $>$ Paper})
        \item Гипотеза: структура графа помогает сопоставлению
      \end{itemize}

      \vspace{0.3em}
      \structure{\large Как обучали}
      \begin{itemize}
        \item Правильная пара (\texttt{Paper}$\leftrightarrow$\texttt{Paper}):
          цель~--- косинус $\to$ 1{,}0
        \item Случайная пара (\texttt{Paper}$\leftrightarrow$\texttt{Reviewer}):
          цель~--- косинус $\to$ 0{,}0
        \item 3 случайных негатива на каждый позитив
        \item 21 фолд leave-one-out: 20 пар учим, одну проверяем
      \end{itemize}

    \column{0.44\textwidth}
      \structure{\large Итог (средний F1 по 21 паре)}

      {\footnotesize
      \begin{tabular}{@{}lc@{\hspace{10pt}}c@{}}
        \toprule
        Подход & F1 & $\Delta$ \\
        \midrule
        MiniLM, жадный 1:1 & 0{,}156 & \\
        GNN (1 слой) & 0{,}204 & $+0{,}048$ \\
        \midrule
        MiniLM, ближайший & 0{,}314 & $-0{,}110$ \\
        \bottomrule
      \end{tabular}}

      \vspace{0.3em}
      \structure{\large Почему не сработало}
      \begin{enumerate}
        \item \textbf{Мало данных}: 20 реальных пар ({\textasciitilde}250 позитивов) + синтетика из тех же онтологий.
        \item \textbf{Структура не переносится}: \texttt{Paper} под \texttt{Document} здесь и под \texttt{Contribution} там.
        \item \textbf{Схема отбора пар у GNN хуже}: тот же матчер даёт $0{,}156 \to 0{,}204$, но выбор ближайшего узла~--- больше.
      \end{enumerate}
  \end{columns}
\end{frame}

% --- 9. Итоговая таблица подходов ---
\begin{frame}{Итоговая таблица (6 <<сложных>> пар, лучший режим LLM)}
  \centering
  \footnotesize

  \begin{tabular}{@{}lcccccc@{}}
    \toprule
    \textbf{Пара} & \textbf{LaBSE} & \textbf{MiniLM} & \textbf{StrEq} & \textbf{GNN} & \textbf{Лучший LLM} & \textbf{LLM F1} \\
    \midrule
    cmt $\leftrightarrow$ confOf
      & 0{,}308 & 0{,}308 & 0{,}533 & 0{,}343 & DeepSeek CP & 0{,}750 \\
    cmt $\leftrightarrow$ conference
      & 0{,}450 & 0{,}450 & 0{,}375 & 0{,}108 & GPT C & 0{,}640 \\
    confOf $\leftrightarrow$ conference
      & 0{,}367 & 0{,}367 & 0{,}737 & 0{,}293 & DS C / BM25 CP & 0{,}762 \\
    cmt $\leftrightarrow$ edas
      & 0{,}432 & 0{,}432 & \textbf{1{,}000} & 0{,}111 & BM25 CC & \textbf{1{,}000} \\
    confOf $\leftrightarrow$ ekaw
      & 0{,}586 & 0{,}586 & 0{,}552 & 0{,}154 & BM25 C & 0{,}667 \\
    conference $\leftrightarrow$ ekaw
      & 0{,}439 & 0{,}463 & 0{,}457 & 0{,}209 & GPT CC & 0{,}577 \\
    \bottomrule
  \end{tabular}

  \vspace{1em}
  \begin{itemize}
    \item LLM превосходит все не-LLM подходы на всех парах, кроме cmt$\leftrightarrow$edas (потолок)
    \item Embedding (LaBSE, MiniLM) лучше StringEquiv на парах таксономий с различной лексикой
    \item GNN систематически хуже raw эмбеддингов
    \vspace{0.5em}
    \item LLM с любым ретривером (MiniLM, BM25, Hybrid) дают лучшие результаты на парах с различной лексикой
    \item StringEquiv -- неожиданно сильный бейзлайн для данного датасета, что подозрительно
    \item GNN не удалось оценить по достоинству из-за недостатка данных
  \end{itemize}

\end{frame}

% --- 10. Выводы ---
\begin{frame}{Выводы}
  \begin{enumerate}
    \item \structure{LLM + эмбеддинги (LLMs4OM) ~--- неплохое решение, если у вас много ресурсов и времени}

    \vspace{0.4em}
    \item \structure{Данных OAEI (Conference Track) недостаточно для развития темы}

    \item \structure{\texttt{name\_only}~--- лучше работает с BERTами: F1 в 1,5 раза выше, чем у имени вместе с путём и атрибутами}

    \vspace{0.4em}
    \item \structure{Графовые сети не дают выигрыша при недостатке данных}

    \vspace{0.4em}
    \item \structure{Требуется дальнейшее исследование возможностей гибридного поиска в задаче OM}

      {\small BM25 лучший на 8/21 парах, Hybrid~--- на 5/21; в среднем оба лучше чистого
      MiniLM-ретривера (0{,}662 и 0{,}685 против 0{,}644).}
  \end{enumerate}
\end{frame}

% --- 11. Направления для дальнейшей работы ---
\begin{frame}{Направления для дальнейшей работы}
  \begin{itemize}
    \item \structure{Проверка согласованности таксономий} — вторая половина темы:
      reasoner (HermiT, Pellet) для OWL + LLM-проверка подозрительных связей

    \vspace{0.3em}
    \item \structure{Пробная интеграция лучших решений в проприетарную ГИС}

    \vspace{0.3em}
    \item \structure{Бо́льший датасет}: синтетика через LLM, полный OntoFarm —
      на нём GNN-архитектуры (GraphSAGE, GAT) и масштабирование LLM через blocking

    \vspace{0.3em}
    \item \structure{Анализ ошибок LLM}: систематические паттерны, ансамблирование режимов
  \end{itemize}
\end{frame}

% --- 12. Спасибо ---
{
\setbeamertemplate{footline}{}
\begin{frame}[plain]
  \begin{center}
    \vspace{2cm}
    {\usebeamercolor[fg]{structure}\fontsize{36}{42}\selectfont Спасибо за внимание!}\\[1cm]
  \end{center}
\end{frame}
}

\end{document}
